\mxfront{terminale-e2}
\mxnewpage\phantomsection\label{sujet:2-1}\addcontentsline{toc}{chapter}{DS2 · Zéro 1}

\section*{DS2 · Zéro 1}

\mxbanner{SUJET ZÉRO · ENTRAÎNEMENT}

\textbf{Durée : 2 h · Note sur 20 · \mxTallied{ex-2-1} exercices indépendants.}\\Calculatrice autorisée. Toute réponse doit être justifiée ; les démarches incomplètes seront prises en compte. Donner les valeurs exactes ; arrondir les probabilités à $10^{-6}$ et les autres approximations à $10^{-4}$, sauf précision différente.

\smallskip\textbf{Nom :}\enspace\dotfill\quad\textbf{Classe :}\enspace\dotfill\quad\textbf{Date :}\enspace\dotfill

\mxtally{ex-2-1}
\subsection*{Exercice 1 · Limites calculées et justifiées \hfill 5 points}Pour tout entier $n\ge1$, on définit $u_n=\frac{3n^2+4n-1}{4n^2+3}$ et $v_n=\sqrt{n^2+4n}-n$.

\begin{enumerate}

\item Déterminer la limite de $(u_n)$ en détaillant la transformation de son expression. \hfill\textit{(1 pt)}

\item Montrer que, pour tout entier $n\ge1$, $v_n=\frac{4}{\sqrt{1+\frac{4}n}+1}$. \hfill\textit{(1 pt)}

\item En déduire la limite de $(v_n)$. \hfill\textit{(1 pt)}

\item On définit $w_n=3n^2-4n$ pour $n\ge1$. Déterminer la limite de $(w_n)$. \hfill\textit{(1 pt)}

\item Un élève affirme que la différence de deux suites tendant vers $+\infty$ tend toujours vers $0$. Cette affirmation est-elle vraie ? Justifier. \hfill\textit{(1 pt)}

\end{enumerate}

\mxtally{ex-2-1}
\subsection*{Exercice 2 · Vrai ou faux : raisonnements sur les limites \hfill 5 points}Pour chacune des affirmations suivantes, indiquer si elle est vraie ou fausse et justifier la réponse. Une réponse non justifiée ne rapporte aucun point. Les limites des suites sont prises lorsque $n$ tend vers $+\infty$.

\begin{enumerate}

\item Si $u_n\to 2$ et $v_n\to-5$, alors $\frac{u_n}{v_n}\to-\frac{2}{5}$. \hfill\textit{(1 pt)}

\item La suite $u_n=2+(-1)^n$ converge vers $2$. \hfill\textit{(1 pt)}

\item Si $u_n\to0$ et $v_n\to+\infty$, alors $u_nv_n\to0$. \hfill\textit{(1 pt)}

\item La suite $w_n=\frac{\sqrt{n^2+3n}}{n}$ pour $n\ge1$ converge vers $1$. \hfill\textit{(1 pt)}

\item Pour $u_n=\frac{2n+5}{n}$, on a $|u_n-2|<\num{0,01}$ dès que $n> 500$. \hfill\textit{(1 pt)}

\end{enumerate}

\mxnewpage\mxcontinued{DS2 · Zéro 1 · suite}

\mxtally{ex-2-1}
\subsection*{Exercice 3 · Calcul vectoriel dans une base \hfill 5 points}Dans une base $\pt{\vi;\vj;\vec k}$ de l'espace, on définit les vecteurs
\[\begin{aligned}\vec u&=\vi+\vec k,&\vec v&=\vj+\vec k,\\
\vec w&=2\vi+6\vj+8\vec k,&\vec z&=2\vi+6\vj+9\vec k.\end{aligned}\]

\begin{enumerate}

\item Justifier que $\vec u$ et $\vec v$ ne sont pas colinéaires. \hfill\textit{(1 pt)}

\item Écrire $\vec w$ comme combinaison linéaire de $\vec u$ et $\vec v$. \hfill\textit{(\num{0,5} pt)}

\item Les vecteurs $\vec u$, $\vec v$ et $\vec z$ sont-ils coplanaires ? Justifier. \hfill\textit{(\num{1,5} pt)}

\item Exprimer le vecteur $6\vec u-2\vec v$ dans la base $\pt{\vi;\vj;\vec k}$. \hfill\textit{(1 pt)}

\item On choisit un point $A$, puis on définit $B$, $C$ et $D$ par $\vect{AB}=\vec u$, $\vect{AC}=\vec v$ et $\vect{AD}=\vec w$. Démontrer que $A$, $B$, $C$ et $D$ sont coplanaires. \hfill\textit{(1 pt)}

\end{enumerate}

\mxtally{ex-2-1}
\subsection*{Exercice 4 · Géométrie d'un parallélépipède \hfill 5 points}On considère un parallélépipède $ABCDEFGH$ non aplati. Les sommets $E$, $F$, $G$ et $H$ sont les images respectives de $A$, $B$, $C$ et $D$ par la translation de vecteur $\vect{AE}$.
On pose $\vi=\vect{AB}$, $\vj=\vect{AD}$ et $\vec k=\vect{AE}$.
Les points $I$, $J$ et $K$ sont définis par
\[\vect{AI}=\frac1{18}\vi,\qquad\vect{AJ}=\frac1{18}\vj,\qquad\vect{AK}=\frac1{18}\vec k.\]

\begin{enumerate}

\item Exprimer $\vect{AG}$ et $\vect{BH}$ dans la base $\pt{\vi;\vj;\vec k}$. \hfill\textit{(1 pt)}

\item Démontrer que les droites $(IJ)$ et $(BD)$ sont parallèles. \hfill\textit{(1 pt)}

\item Étudier la position relative des droites $(AB)$ et $(CG)$ : sont-elles parallèles, sécantes ou non coplanaires ? Justifier. \hfill\textit{(1 pt)}

\item On note $O$ le centre du parallélépipède. Le point $O$ appartient-il au plan $(ABG)$ ? Justifier. \hfill\textit{(1 pt)}

\item Démontrer que les points $I$, $J$ et $K$ ne sont pas alignés. \hfill\textit{(1 pt)}

\end{enumerate}

\ifmxteacher\mxnewpage

\section*{Corrigé · DS2 · Zéro 1}\label{cor:2-1}

\subsection*{Exercice 1 · Limites calculées et justifiées}\begin{enumerate}

\item \textit{[1 pt]} Pour $n\ge1$, on divise par $n^2\ne0$ :
\[u_n=\frac{3+\frac{4}{n}-\frac{1}{n^2}}{4+\frac{3}{n^2}}.\]
Les termes $\frac{1}{n}$ et $\frac{1}{n^2}$ tendent vers $0$. Le numérateur tend vers $3$ et le dénominateur vers $4\ne0$ ; par quotient, $u_n\to\frac{3}{4}$.

\item \textit{[1 pt]} Le conjugué est positif. Pour $n\ge1$, $\sqrt{n^2}=n$, d'où
\[\begin{aligned}v_n&=\frac{(n^2+4n)-n^2}{\sqrt{n^2+4n}+n}\\&=\frac{4n}{n(\sqrt{1+\frac{4}{n}}+1)}=\frac{4}{\sqrt{1+\frac{4}{n}}+1}.\end{aligned}\]

\item \textit{[1 pt]} Le dénominateur tend vers $2$ ; limite $2$.

\item \textit{[1 pt]} $w_n=n^2(3-\frac{4}{n})\to+\infty$ ; le second facteur tend vers $3>0$.

\item \textit{[1 pt]} Non : $(2n)-n=n\to+\infty$. La forme $\infty-\infty$ est indéterminée : elle ne permet pas de conclure sans transformation.

\end{enumerate}

\subsection*{Exercice 2 · Vrai ou faux : raisonnements sur les limites}\begin{enumerate}

\item \textit{[1 pt]} Vraie : le dénominateur a une limite non nulle ; la limite du quotient est $-\frac{2}{5}$.

\item \textit{[1 pt]} Fausse : les termes pairs valent $3$, les termes impairs $1$. Les deux valeurs restent à distance $1$ de $2$.

\item \textit{[1 pt]} Fausse : $u_n=\frac{1}{n}$, $v_n=n$ pour $n\ge1$ donnent un produit constant égal à $1$.

\item \textit{[1 pt]} Vraie : $w_n=\sqrt{1+\frac{3}{n}}\to1$.

\item \textit{[1 pt]} Vraie : l'écart vaut $\frac{5}{n}$ ; $n>500$ donne $\frac{5}{n}<\num{0,01}$.

\end{enumerate}

\subsection*{Exercice 3 · Calcul vectoriel dans une base}\begin{enumerate}

\item \textit{[1 pt]} Dans une égalité $\vec v=t\vec u$, le coefficient de $\vj $ serait $1=0$, contradiction par unicité des coordonnées dans une base.

\item \textit{[\num{0,5} pt]} En développant, $2\vec u+6\vec v=2(\vi +\vec k)+6(\vj +\vec k)=2\vi +6\vj +8\vec k=\vec w$.

\item \textit{[\num{1,5} pt]} Comme $\vec u,\vec v$ ne sont pas colinéaires, la coplanarité équivaudrait à $\vec z=\lambda\vec u+\mu\vec v$. L'unicité de la décomposition dans la base impose
\[\lambda=2,\qquad\mu=6,\qquad\lambda+\mu=9.\]
Les deux premières égalités donnent $\lambda+\mu=8$, en contradiction avec la troisième. Les trois vecteurs ne sont donc pas coplanaires.

\item \textit{[1 pt]} $6\vec u-2\vec v=6(\vi +\vec k)-2(\vj +\vec k)=6\vi -2\vj +4\vec k$.

\item \textit{[1 pt]} $\vect{AB}$ et $\vect{AC}$ sont non colinéaires et $\vect{AD}$ est leur combinaison linéaire : $D\in(ABC)$.

\end{enumerate}

\subsection*{Exercice 4 · Géométrie d'un parallélépipède}\begin{enumerate}

\item \textit{[1 pt]} Par Chasles, $\vect{AG}=\vect{AB}+\vect{BC}+\vect{CG}=\vi +\vj +\vec k$. De même, $\vect{BH}=\vect{BA}+\vect{AD}+\vect{DH}=-\vi +\vj +\vec k$.

\item \textit{[1 pt]} Par Chasles, $\vect{IJ}=\vect{AJ}-\vect{AI}=\frac{\vj -\vi }{18}=\frac{\vect{BD}}{18}$. Ces vecteurs sont non nuls et colinéaires : $(IJ)$ et $(BD)$ sont parallèles.

\item \textit{[1 pt]} Directeurs $\vi $ et $\vec k$, non colinéaires. Une intersection imposerait $t\vi =\vi +\vj +s\vec k$, impossible par le coefficient de $\vj $. Elles sont non coplanaires.

\item \textit{[1 pt]} $\vect{AO}=\frac{\vect{AG}}{2}$, donc $O\in(AG)\subset(ABG)$.

\item \textit{[1 pt]} Les points vérifient $\vect{IJ}=\frac{-\vi +\vj }{18}$ et $\vect{IK}=\frac{-\vi +\vec k}{18}$. Si $\vect{IK}=t\vect{IJ}$, le coefficient de $\vec k$ donnerait $\frac{1}{18}=0$, impossible. Les vecteurs ne sont pas colinéaires, donc $I,J,K$ ne sont pas alignés.

\end{enumerate}

\fi

\mxnewpage\phantomsection\label{sujet:2-2}\addcontentsline{toc}{chapter}{DS2 · Zéro 2}

\section*{DS2 · Zéro 2}

\mxbanner{SUJET ZÉRO · ENTRAÎNEMENT}

\textbf{Durée : 2 h · Note sur 20 · \mxTallied{ex-2-2} exercices indépendants.}\\Calculatrice autorisée. Toute réponse doit être justifiée ; les démarches incomplètes seront prises en compte. Donner les valeurs exactes ; arrondir les probabilités à $10^{-6}$ et les autres approximations à $10^{-4}$, sauf précision différente.

\smallskip\textbf{Nom :}\enspace\dotfill\quad\textbf{Classe :}\enspace\dotfill\quad\textbf{Date :}\enspace\dotfill

\mxtally{ex-2-2}
\subsection*{Exercice 1 · Limites calculées et justifiées \hfill 5 points}Pour tout entier $n\ge1$, on définit $u_n=\frac{2n^2+6n-1}{3n^2+1}$ et $v_n=\sqrt{n^2+6n}-n$.

\begin{enumerate}

\item Déterminer la limite de $(u_n)$ en détaillant la transformation de son expression. \hfill\textit{(1 pt)}

\item Montrer que, pour tout entier $n\ge1$, $v_n=\frac{6}{\sqrt{1+\frac{6}n}+1}$. \hfill\textit{(1 pt)}

\item En déduire la limite de $(v_n)$. \hfill\textit{(1 pt)}

\item On définit $w_n=2n^2-6n$ pour $n\ge1$. Déterminer la limite de $(w_n)$. \hfill\textit{(1 pt)}

\item Un élève affirme que la différence de deux suites tendant vers $+\infty$ tend toujours vers $0$. Cette affirmation est-elle vraie ? Justifier. \hfill\textit{(1 pt)}

\end{enumerate}

\mxtally{ex-2-2}
\subsection*{Exercice 2 · Vrai ou faux : raisonnements sur les limites \hfill 5 points}Pour chacune des affirmations suivantes, indiquer si elle est vraie ou fausse et justifier la réponse. Une réponse non justifiée ne rapporte aucun point. Les limites des suites sont prises lorsque $n$ tend vers $+\infty$.

\begin{enumerate}

\item Si $u_n\to 5$ et $v_n\to-6$, alors $\frac{u_n}{v_n}\to-\frac{5}{6}$. \hfill\textit{(1 pt)}

\item La suite $u_n=5+(-1)^n$ converge vers $5$. \hfill\textit{(1 pt)}

\item Si $u_n\to0$ et $v_n\to+\infty$, alors $u_nv_n\to0$. \hfill\textit{(1 pt)}

\item La suite $w_n=\frac{\sqrt{n^2+n}}{n}$ pour $n\ge1$ converge vers $1$. \hfill\textit{(1 pt)}

\item Pour $u_n=\frac{5n+6}{n}$, on a $|u_n-5|<\num{0,01}$ dès que $n> 600$. \hfill\textit{(1 pt)}

\end{enumerate}

\mxnewpage\mxcontinued{DS2 · Zéro 2 · suite}

\mxtally{ex-2-2}
\subsection*{Exercice 3 · Calcul vectoriel dans une base \hfill 5 points}Dans une base $\pt{\vi;\vj;\vec k}$ de l'espace, on définit les vecteurs
\[\begin{aligned}\vec u&=\vi+\vec k,&\vec v&=\vj+\vec k,\\
\vec w&=4\vi+7\vj+11\vec k,&\vec z&=4\vi+7\vj+12\vec k.\end{aligned}\]

\begin{enumerate}

\item Justifier que $\vec u$ et $\vec v$ ne sont pas colinéaires. \hfill\textit{(1 pt)}

\item Écrire $\vec w$ comme combinaison linéaire de $\vec u$ et $\vec v$. \hfill\textit{(\num{0,5} pt)}

\item Les vecteurs $\vec u$, $\vec v$ et $\vec z$ sont-ils coplanaires ? Justifier. \hfill\textit{(\num{1,5} pt)}

\item Exprimer le vecteur $7\vec u-4\vec v$ dans la base $\pt{\vi;\vj;\vec k}$. \hfill\textit{(1 pt)}

\item On choisit un point $A$, puis on définit $B$, $C$ et $D$ par $\vect{AB}=\vec u$, $\vect{AC}=\vec v$ et $\vect{AD}=\vec w$. Démontrer que $A$, $B$, $C$ et $D$ sont coplanaires. \hfill\textit{(1 pt)}

\end{enumerate}

\mxtally{ex-2-2}
\subsection*{Exercice 4 · Géométrie d'un parallélépipède \hfill 5 points}On considère un parallélépipède $ABCDEFGH$ non aplati. Les sommets $E$, $F$, $G$ et $H$ sont les images respectives de $A$, $B$, $C$ et $D$ par la translation de vecteur $\vect{AE}$.
On pose $\vi=\vect{AB}$, $\vj=\vect{AD}$ et $\vec k=\vect{AE}$.
Les points $I$, $J$ et $K$ sont définis par
\[\vect{AI}=\frac1{8}\vi,\qquad\vect{AJ}=\frac1{8}\vj,\qquad\vect{AK}=\frac1{8}\vec k.\]

\begin{enumerate}

\item Exprimer $\vect{AG}$ et $\vect{BH}$ dans la base $\pt{\vi;\vj;\vec k}$. \hfill\textit{(1 pt)}

\item Démontrer que les droites $(IJ)$ et $(BD)$ sont parallèles. \hfill\textit{(1 pt)}

\item Étudier la position relative des droites $(AB)$ et $(CG)$ : sont-elles parallèles, sécantes ou non coplanaires ? Justifier. \hfill\textit{(1 pt)}

\item On note $O$ le centre du parallélépipède. Le point $O$ appartient-il au plan $(ABG)$ ? Justifier. \hfill\textit{(1 pt)}

\item Démontrer que les points $I$, $J$ et $K$ ne sont pas alignés. \hfill\textit{(1 pt)}

\end{enumerate}

\ifmxteacher\mxnewpage

\section*{Corrigé · DS2 · Zéro 2}\label{cor:2-2}

\subsection*{Exercice 1 · Limites calculées et justifiées}\begin{enumerate}

\item \textit{[1 pt]} Pour $n\ge1$, on divise par $n^2\ne0$ :
\[u_n=\frac{2+\frac{6}{n}-\frac{1}{n^2}}{3+\frac{1}{n^2}}.\]
Les termes $\frac{1}{n}$ et $\frac{1}{n^2}$ tendent vers $0$. Le numérateur tend vers $2$ et le dénominateur vers $3\ne0$ ; par quotient, $u_n\to\frac{2}{3}$.

\item \textit{[1 pt]} Le conjugué est positif. Pour $n\ge1$, $\sqrt{n^2}=n$, d'où
\[\begin{aligned}v_n&=\frac{(n^2+6n)-n^2}{\sqrt{n^2+6n}+n}\\&=\frac{6n}{n(\sqrt{1+\frac{6}{n}}+1)}=\frac{6}{\sqrt{1+\frac{6}{n}}+1}.\end{aligned}\]

\item \textit{[1 pt]} Le dénominateur tend vers $2$ ; limite $3$.

\item \textit{[1 pt]} $w_n=n^2(2-\frac{6}{n})\to+\infty$ ; le second facteur tend vers $2>0$.

\item \textit{[1 pt]} Non : $(2n)-n=n\to+\infty$. La forme $\infty-\infty$ est indéterminée : elle ne permet pas de conclure sans transformation.

\end{enumerate}

\subsection*{Exercice 2 · Vrai ou faux : raisonnements sur les limites}\begin{enumerate}

\item \textit{[1 pt]} Vraie : le dénominateur a une limite non nulle ; la limite du quotient est $-\frac{5}{6}$.

\item \textit{[1 pt]} Fausse : les termes pairs valent $6$, les termes impairs $4$. Les deux valeurs restent à distance $1$ de $5$.

\item \textit{[1 pt]} Fausse : $u_n=\frac{1}{n}$, $v_n=n$ pour $n\ge1$ donnent un produit constant égal à $1$.

\item \textit{[1 pt]} Vraie : $w_n=\sqrt{1+\frac{1}{n}}\to1$.

\item \textit{[1 pt]} Vraie : l'écart vaut $\frac{6}{n}$ ; $n>600$ donne $\frac{6}{n}<\num{0,01}$.

\end{enumerate}

\subsection*{Exercice 3 · Calcul vectoriel dans une base}\begin{enumerate}

\item \textit{[1 pt]} Dans une égalité $\vec v=t\vec u$, le coefficient de $\vj $ serait $1=0$, contradiction par unicité des coordonnées dans une base.

\item \textit{[\num{0,5} pt]} En développant, $4\vec u+7\vec v=4(\vi +\vec k)+7(\vj +\vec k)=4\vi +7\vj +11\vec k=\vec w$.

\item \textit{[\num{1,5} pt]} Comme $\vec u,\vec v$ ne sont pas colinéaires, la coplanarité équivaudrait à $\vec z=\lambda\vec u+\mu\vec v$. L'unicité de la décomposition dans la base impose
\[\lambda=4,\qquad\mu=7,\qquad\lambda+\mu=12.\]
Les deux premières égalités donnent $\lambda+\mu=11$, en contradiction avec la troisième. Les trois vecteurs ne sont donc pas coplanaires.

\item \textit{[1 pt]} $7\vec u-4\vec v=7(\vi +\vec k)-4(\vj +\vec k)=7\vi -4\vj +3\vec k$.

\item \textit{[1 pt]} $\vect{AB}$ et $\vect{AC}$ sont non colinéaires et $\vect{AD}$ est leur combinaison linéaire : $D\in(ABC)$.

\end{enumerate}

\subsection*{Exercice 4 · Géométrie d'un parallélépipède}\begin{enumerate}

\item \textit{[1 pt]} Par Chasles, $\vect{AG}=\vect{AB}+\vect{BC}+\vect{CG}=\vi +\vj +\vec k$. De même, $\vect{BH}=\vect{BA}+\vect{AD}+\vect{DH}=-\vi +\vj +\vec k$.

\item \textit{[1 pt]} Par Chasles, $\vect{IJ}=\vect{AJ}-\vect{AI}=\frac{\vj -\vi }{8}=\frac{\vect{BD}}{8}$. Ces vecteurs sont non nuls et colinéaires : $(IJ)$ et $(BD)$ sont parallèles.

\item \textit{[1 pt]} Directeurs $\vi $ et $\vec k$, non colinéaires. Une intersection imposerait $t\vi =\vi +\vj +s\vec k$, impossible par le coefficient de $\vj $. Elles sont non coplanaires.

\item \textit{[1 pt]} $\vect{AO}=\frac{\vect{AG}}{2}$, donc $O\in(AG)\subset(ABG)$.

\item \textit{[1 pt]} Les points vérifient $\vect{IJ}=\frac{-\vi +\vj }{8}$ et $\vect{IK}=\frac{-\vi +\vec k}{8}$. Si $\vect{IK}=t\vect{IJ}$, le coefficient de $\vec k$ donnerait $\frac{1}{8}=0$, impossible. Les vecteurs ne sont pas colinéaires, donc $I,J,K$ ne sont pas alignés.

\end{enumerate}

\fi

\ifmxteacher

\mxnewpage\phantomsection\label{sujet:2-3}\addcontentsline{toc}{chapter}{DS2 · A}

\section*{DS2 · A}

\mxbanner*{SUJET RÉSERVÉ À L'ÉVALUATION}

\textbf{Durée : 2 h · Note sur 20 · \mxTallied{ex-2-3} exercices indépendants.}\\Calculatrice autorisée. Toute réponse doit être justifiée ; les démarches incomplètes seront prises en compte. Donner les valeurs exactes ; arrondir les probabilités à $10^{-6}$ et les autres approximations à $10^{-4}$, sauf précision différente.

\smallskip\textbf{Nom :}\enspace\dotfill\quad\textbf{Classe :}\enspace\dotfill\quad\textbf{Date :}\enspace\dotfill

\mxtally{ex-2-3}
\subsection*{Exercice 1 · Limites calculées et justifiées \hfill 5 points}Pour tout entier $n\ge1$, on définit $u_n=\frac{5n^2+7n-1}{6n^2+2}$ et $v_n=\sqrt{n^2+7n}-n$.

\begin{enumerate}

\item Déterminer la limite de $(u_n)$ en détaillant la transformation de son expression. \hfill\textit{(1 pt)}

\item Montrer que, pour tout entier $n\ge1$, $v_n=\frac{7}{\sqrt{1+\frac{7}n}+1}$. \hfill\textit{(1 pt)}

\item En déduire la limite de $(v_n)$. \hfill\textit{(1 pt)}

\item On définit $w_n=5n^2-7n$ pour $n\ge1$. Déterminer la limite de $(w_n)$. \hfill\textit{(1 pt)}

\item Un élève affirme que la différence de deux suites tendant vers $+\infty$ tend toujours vers $0$. Cette affirmation est-elle vraie ? Justifier. \hfill\textit{(1 pt)}

\end{enumerate}

\mxtally{ex-2-3}
\subsection*{Exercice 2 · Vrai ou faux : raisonnements sur les limites \hfill 5 points}Pour chacune des affirmations suivantes, indiquer si elle est vraie ou fausse et justifier la réponse. Une réponse non justifiée ne rapporte aucun point. Les limites des suites sont prises lorsque $n$ tend vers $+\infty$.

\begin{enumerate}

\item Si $u_n\to 4$ et $v_n\to-8$, alors $\frac{u_n}{v_n}\to-\frac{1}{2}$. \hfill\textit{(1 pt)}

\item La suite $u_n=4+(-1)^n$ converge vers $4$. \hfill\textit{(1 pt)}

\item Si $u_n\to0$ et $v_n\to+\infty$, alors $u_nv_n\to0$. \hfill\textit{(1 pt)}

\item La suite $w_n=\frac{\sqrt{n^2+2n}}{n}$ pour $n\ge1$ converge vers $1$. \hfill\textit{(1 pt)}

\item Pour $u_n=\frac{4n+8}{n}$, on a $|u_n-4|<\num{0,01}$ dès que $n> 800$. \hfill\textit{(1 pt)}

\end{enumerate}

\mxnewpage\mxcontinued{DS2 · A · suite}

\mxtally{ex-2-3}
\subsection*{Exercice 3 · Calcul vectoriel dans une base \hfill 5 points}Dans une base $\pt{\vi;\vj;\vec k}$ de l'espace, on définit les vecteurs
\[\begin{aligned}\vec u&=\vi+\vec k,&\vec v&=\vj+\vec k,\\
\vec w&=3\vi+4\vj+7\vec k,&\vec z&=3\vi+4\vj+8\vec k.\end{aligned}\]

\begin{enumerate}

\item Justifier que $\vec u$ et $\vec v$ ne sont pas colinéaires. \hfill\textit{(1 pt)}

\item Écrire $\vec w$ comme combinaison linéaire de $\vec u$ et $\vec v$. \hfill\textit{(\num{0,5} pt)}

\item Les vecteurs $\vec u$, $\vec v$ et $\vec z$ sont-ils coplanaires ? Justifier. \hfill\textit{(\num{1,5} pt)}

\item Exprimer le vecteur $4\vec u-3\vec v$ dans la base $\pt{\vi;\vj;\vec k}$. \hfill\textit{(1 pt)}

\item On choisit un point $A$, puis on définit $B$, $C$ et $D$ par $\vect{AB}=\vec u$, $\vect{AC}=\vec v$ et $\vect{AD}=\vec w$. Démontrer que $A$, $B$, $C$ et $D$ sont coplanaires. \hfill\textit{(1 pt)}

\end{enumerate}

\mxtally{ex-2-3}
\subsection*{Exercice 4 · Géométrie d'un parallélépipède \hfill 5 points}On considère un parallélépipède $ABCDEFGH$ non aplati. Les sommets $E$, $F$, $G$ et $H$ sont les images respectives de $A$, $B$, $C$ et $D$ par la translation de vecteur $\vect{AE}$.
On pose $\vi=\vect{AB}$, $\vj=\vect{AD}$ et $\vec k=\vect{AE}$.
Les points $I$, $J$ et $K$ sont définis par
\[\vect{AI}=\frac1{15}\vi,\qquad\vect{AJ}=\frac1{15}\vj,\qquad\vect{AK}=\frac1{15}\vec k.\]

\begin{enumerate}

\item Exprimer $\vect{AG}$ et $\vect{BH}$ dans la base $\pt{\vi;\vj;\vec k}$. \hfill\textit{(1 pt)}

\item Démontrer que les droites $(IJ)$ et $(BD)$ sont parallèles. \hfill\textit{(1 pt)}

\item Étudier la position relative des droites $(AB)$ et $(CG)$ : sont-elles parallèles, sécantes ou non coplanaires ? Justifier. \hfill\textit{(1 pt)}

\item On note $O$ le centre du parallélépipède. Le point $O$ appartient-il au plan $(ABG)$ ? Justifier. \hfill\textit{(1 pt)}

\item Démontrer que les points $I$, $J$ et $K$ ne sont pas alignés. \hfill\textit{(1 pt)}

\end{enumerate}

\ifmxteacher\mxnewpage

\section*{Corrigé · DS2 · A}\label{cor:2-3}

\subsection*{Exercice 1 · Limites calculées et justifiées}\begin{enumerate}

\item \textit{[1 pt]} Pour $n\ge1$, on divise par $n^2\ne0$ :
\[u_n=\frac{5+\frac{7}{n}-\frac{1}{n^2}}{6+\frac{2}{n^2}}.\]
Les termes $\frac{1}{n}$ et $\frac{1}{n^2}$ tendent vers $0$. Le numérateur tend vers $5$ et le dénominateur vers $6\ne0$ ; par quotient, $u_n\to\frac{5}{6}$.

\item \textit{[1 pt]} Le conjugué est positif. Pour $n\ge1$, $\sqrt{n^2}=n$, d'où
\[\begin{aligned}v_n&=\frac{(n^2+7n)-n^2}{\sqrt{n^2+7n}+n}\\&=\frac{7n}{n(\sqrt{1+\frac{7}{n}}+1)}=\frac{7}{\sqrt{1+\frac{7}{n}}+1}.\end{aligned}\]

\item \textit{[1 pt]} Le dénominateur tend vers $2$ ; limite $\frac{7}{2}$.

\item \textit{[1 pt]} $w_n=n^2(5-\frac{7}{n})\to+\infty$ ; le second facteur tend vers $5>0$.

\item \textit{[1 pt]} Non : $(2n)-n=n\to+\infty$. La forme $\infty-\infty$ est indéterminée : elle ne permet pas de conclure sans transformation.

\end{enumerate}

\subsection*{Exercice 2 · Vrai ou faux : raisonnements sur les limites}\begin{enumerate}

\item \textit{[1 pt]} Vraie : le dénominateur a une limite non nulle ; la limite du quotient est $-\frac{1}{2}$.

\item \textit{[1 pt]} Fausse : les termes pairs valent $5$, les termes impairs $3$. Les deux valeurs restent à distance $1$ de $4$.

\item \textit{[1 pt]} Fausse : $u_n=\frac{1}{n}$, $v_n=n$ pour $n\ge1$ donnent un produit constant égal à $1$.

\item \textit{[1 pt]} Vraie : $w_n=\sqrt{1+\frac{2}{n}}\to1$.

\item \textit{[1 pt]} Vraie : l'écart vaut $\frac{8}{n}$ ; $n>800$ donne $\frac{8}{n}<\num{0,01}$.

\end{enumerate}

\subsection*{Exercice 3 · Calcul vectoriel dans une base}\begin{enumerate}

\item \textit{[1 pt]} Dans une égalité $\vec v=t\vec u$, le coefficient de $\vj $ serait $1=0$, contradiction par unicité des coordonnées dans une base.

\item \textit{[\num{0,5} pt]} En développant, $3\vec u+4\vec v=3(\vi +\vec k)+4(\vj +\vec k)=3\vi +4\vj +7\vec k=\vec w$.

\item \textit{[\num{1,5} pt]} Comme $\vec u,\vec v$ ne sont pas colinéaires, la coplanarité équivaudrait à $\vec z=\lambda\vec u+\mu\vec v$. L'unicité de la décomposition dans la base impose
\[\lambda=3,\qquad\mu=4,\qquad\lambda+\mu=8.\]
Les deux premières égalités donnent $\lambda+\mu=7$, en contradiction avec la troisième. Les trois vecteurs ne sont donc pas coplanaires.

\item \textit{[1 pt]} $4\vec u-3\vec v=4(\vi +\vec k)-3(\vj +\vec k)=4\vi -3\vj +\vec k$.

\item \textit{[1 pt]} $\vect{AB}$ et $\vect{AC}$ sont non colinéaires et $\vect{AD}$ est leur combinaison linéaire : $D\in(ABC)$.

\end{enumerate}

\subsection*{Exercice 4 · Géométrie d'un parallélépipède}\begin{enumerate}

\item \textit{[1 pt]} Par Chasles, $\vect{AG}=\vect{AB}+\vect{BC}+\vect{CG}=\vi +\vj +\vec k$. De même, $\vect{BH}=\vect{BA}+\vect{AD}+\vect{DH}=-\vi +\vj +\vec k$.

\item \textit{[1 pt]} Par Chasles, $\vect{IJ}=\vect{AJ}-\vect{AI}=\frac{\vj -\vi }{15}=\frac{\vect{BD}}{15}$. Ces vecteurs sont non nuls et colinéaires : $(IJ)$ et $(BD)$ sont parallèles.

\item \textit{[1 pt]} Directeurs $\vi $ et $\vec k$, non colinéaires. Une intersection imposerait $t\vi =\vi +\vj +s\vec k$, impossible par le coefficient de $\vj $. Elles sont non coplanaires.

\item \textit{[1 pt]} $\vect{AO}=\frac{\vect{AG}}{2}$, donc $O\in(AG)\subset(ABG)$.

\item \textit{[1 pt]} Les points vérifient $\vect{IJ}=\frac{-\vi +\vj }{15}$ et $\vect{IK}=\frac{-\vi +\vec k}{15}$. Si $\vect{IK}=t\vect{IJ}$, le coefficient de $\vec k$ donnerait $\frac{1}{15}=0$, impossible. Les vecteurs ne sont pas colinéaires, donc $I,J,K$ ne sont pas alignés.

\end{enumerate}

\fi

\fi

\ifmxteacher

\mxnewpage\phantomsection\label{sujet:2-4}\addcontentsline{toc}{chapter}{DS2 · B}

\section*{DS2 · B}

\mxbanner*{SUJET RÉSERVÉ À L'ÉVALUATION}

\textbf{Durée : 2 h · Note sur 20 · \mxTallied{ex-2-4} exercices indépendants.}\\Calculatrice autorisée. Toute réponse doit être justifiée ; les démarches incomplètes seront prises en compte. Donner les valeurs exactes ; arrondir les probabilités à $10^{-6}$ et les autres approximations à $10^{-4}$, sauf précision différente.

\smallskip\textbf{Nom :}\enspace\dotfill\quad\textbf{Classe :}\enspace\dotfill\quad\textbf{Date :}\enspace\dotfill

\mxtally{ex-2-4}
\subsection*{Exercice 1 · Limites calculées et justifiées \hfill 5 points}Pour tout entier $n\ge1$, on définit $u_n=\frac{4n^2+4n-1}{5n^2+3}$ et $v_n=\sqrt{n^2+4n}-n$.

\begin{enumerate}

\item Déterminer la limite de $(u_n)$ en détaillant la transformation de son expression. \hfill\textit{(1 pt)}

\item Montrer que, pour tout entier $n\ge1$, $v_n=\frac{4}{\sqrt{1+\frac{4}n}+1}$. \hfill\textit{(1 pt)}

\item En déduire la limite de $(v_n)$. \hfill\textit{(1 pt)}

\item On définit $w_n=4n^2-4n$ pour $n\ge1$. Déterminer la limite de $(w_n)$. \hfill\textit{(1 pt)}

\item Un élève affirme que la différence de deux suites tendant vers $+\infty$ tend toujours vers $0$. Cette affirmation est-elle vraie ? Justifier. \hfill\textit{(1 pt)}

\end{enumerate}

\mxtally{ex-2-4}
\subsection*{Exercice 2 · Vrai ou faux : raisonnements sur les limites \hfill 5 points}Pour chacune des affirmations suivantes, indiquer si elle est vraie ou fausse et justifier la réponse. Une réponse non justifiée ne rapporte aucun point. Les limites des suites sont prises lorsque $n$ tend vers $+\infty$.

\begin{enumerate}

\item Si $u_n\to 3$ et $v_n\to-5$, alors $\frac{u_n}{v_n}\to-\frac{3}{5}$. \hfill\textit{(1 pt)}

\item La suite $u_n=3+(-1)^n$ converge vers $3$. \hfill\textit{(1 pt)}

\item Si $u_n\to0$ et $v_n\to+\infty$, alors $u_nv_n\to0$. \hfill\textit{(1 pt)}

\item La suite $w_n=\frac{\sqrt{n^2+3n}}{n}$ pour $n\ge1$ converge vers $1$. \hfill\textit{(1 pt)}

\item Pour $u_n=\frac{3n+5}{n}$, on a $|u_n-3|<\num{0,01}$ dès que $n> 500$. \hfill\textit{(1 pt)}

\end{enumerate}

\mxnewpage\mxcontinued{DS2 · B · suite}

\mxtally{ex-2-4}
\subsection*{Exercice 3 · Calcul vectoriel dans une base \hfill 5 points}Dans une base $\pt{\vi;\vj;\vec k}$ de l'espace, on définit les vecteurs
\[\begin{aligned}\vec u&=\vi+\vec k,&\vec v&=\vj+\vec k,\\
\vec w&=3\vi+6\vj+9\vec k,&\vec z&=3\vi+6\vj+10\vec k.\end{aligned}\]

\begin{enumerate}

\item Justifier que $\vec u$ et $\vec v$ ne sont pas colinéaires. \hfill\textit{(1 pt)}

\item Écrire $\vec w$ comme combinaison linéaire de $\vec u$ et $\vec v$. \hfill\textit{(\num{0,5} pt)}

\item Les vecteurs $\vec u$, $\vec v$ et $\vec z$ sont-ils coplanaires ? Justifier. \hfill\textit{(\num{1,5} pt)}

\item Exprimer le vecteur $6\vec u-3\vec v$ dans la base $\pt{\vi;\vj;\vec k}$. \hfill\textit{(1 pt)}

\item On choisit un point $A$, puis on définit $B$, $C$ et $D$ par $\vect{AB}=\vec u$, $\vect{AC}=\vec v$ et $\vect{AD}=\vec w$. Démontrer que $A$, $B$, $C$ et $D$ sont coplanaires. \hfill\textit{(1 pt)}

\end{enumerate}

\mxtally{ex-2-4}
\subsection*{Exercice 4 · Géométrie d'un parallélépipède \hfill 5 points}On considère un parallélépipède $ABCDEFGH$ non aplati. Les sommets $E$, $F$, $G$ et $H$ sont les images respectives de $A$, $B$, $C$ et $D$ par la translation de vecteur $\vect{AE}$.
On pose $\vi=\vect{AB}$, $\vj=\vect{AD}$ et $\vec k=\vect{AE}$.
Les points $I$, $J$ et $K$ sont définis par
\[\vect{AI}=\frac1{5}\vi,\qquad\vect{AJ}=\frac1{5}\vj,\qquad\vect{AK}=\frac1{5}\vec k.\]

\begin{enumerate}

\item Exprimer $\vect{AG}$ et $\vect{BH}$ dans la base $\pt{\vi;\vj;\vec k}$. \hfill\textit{(1 pt)}

\item Démontrer que les droites $(IJ)$ et $(BD)$ sont parallèles. \hfill\textit{(1 pt)}

\item Étudier la position relative des droites $(AB)$ et $(CG)$ : sont-elles parallèles, sécantes ou non coplanaires ? Justifier. \hfill\textit{(1 pt)}

\item On note $O$ le centre du parallélépipède. Le point $O$ appartient-il au plan $(ABG)$ ? Justifier. \hfill\textit{(1 pt)}

\item Démontrer que les points $I$, $J$ et $K$ ne sont pas alignés. \hfill\textit{(1 pt)}

\end{enumerate}

\ifmxteacher\mxnewpage

\section*{Corrigé · DS2 · B}\label{cor:2-4}

\subsection*{Exercice 1 · Limites calculées et justifiées}\begin{enumerate}

\item \textit{[1 pt]} Pour $n\ge1$, on divise par $n^2\ne0$ :
\[u_n=\frac{4+\frac{4}{n}-\frac{1}{n^2}}{5+\frac{3}{n^2}}.\]
Les termes $\frac{1}{n}$ et $\frac{1}{n^2}$ tendent vers $0$. Le numérateur tend vers $4$ et le dénominateur vers $5\ne0$ ; par quotient, $u_n\to\frac{4}{5}$.

\item \textit{[1 pt]} Le conjugué est positif. Pour $n\ge1$, $\sqrt{n^2}=n$, d'où
\[\begin{aligned}v_n&=\frac{(n^2+4n)-n^2}{\sqrt{n^2+4n}+n}\\&=\frac{4n}{n(\sqrt{1+\frac{4}{n}}+1)}=\frac{4}{\sqrt{1+\frac{4}{n}}+1}.\end{aligned}\]

\item \textit{[1 pt]} Le dénominateur tend vers $2$ ; limite $2$.

\item \textit{[1 pt]} $w_n=n^2(4-\frac{4}{n})\to+\infty$ ; le second facteur tend vers $4>0$.

\item \textit{[1 pt]} Non : $(2n)-n=n\to+\infty$. La forme $\infty-\infty$ est indéterminée : elle ne permet pas de conclure sans transformation.

\end{enumerate}

\subsection*{Exercice 2 · Vrai ou faux : raisonnements sur les limites}\begin{enumerate}

\item \textit{[1 pt]} Vraie : le dénominateur a une limite non nulle ; la limite du quotient est $-\frac{3}{5}$.

\item \textit{[1 pt]} Fausse : les termes pairs valent $4$, les termes impairs $2$. Les deux valeurs restent à distance $1$ de $3$.

\item \textit{[1 pt]} Fausse : $u_n=\frac{1}{n}$, $v_n=n$ pour $n\ge1$ donnent un produit constant égal à $1$.

\item \textit{[1 pt]} Vraie : $w_n=\sqrt{1+\frac{3}{n}}\to1$.

\item \textit{[1 pt]} Vraie : l'écart vaut $\frac{5}{n}$ ; $n>500$ donne $\frac{5}{n}<\num{0,01}$.

\end{enumerate}

\subsection*{Exercice 3 · Calcul vectoriel dans une base}\begin{enumerate}

\item \textit{[1 pt]} Dans une égalité $\vec v=t\vec u$, le coefficient de $\vj $ serait $1=0$, contradiction par unicité des coordonnées dans une base.

\item \textit{[\num{0,5} pt]} En développant, $3\vec u+6\vec v=3(\vi +\vec k)+6(\vj +\vec k)=3\vi +6\vj +9\vec k=\vec w$.

\item \textit{[\num{1,5} pt]} Comme $\vec u,\vec v$ ne sont pas colinéaires, la coplanarité équivaudrait à $\vec z=\lambda\vec u+\mu\vec v$. L'unicité de la décomposition dans la base impose
\[\lambda=3,\qquad\mu=6,\qquad\lambda+\mu=10.\]
Les deux premières égalités donnent $\lambda+\mu=9$, en contradiction avec la troisième. Les trois vecteurs ne sont donc pas coplanaires.

\item \textit{[1 pt]} $6\vec u-3\vec v=6(\vi +\vec k)-3(\vj +\vec k)=6\vi -3\vj +3\vec k$.

\item \textit{[1 pt]} $\vect{AB}$ et $\vect{AC}$ sont non colinéaires et $\vect{AD}$ est leur combinaison linéaire : $D\in(ABC)$.

\end{enumerate}

\subsection*{Exercice 4 · Géométrie d'un parallélépipède}\begin{enumerate}

\item \textit{[1 pt]} Par Chasles, $\vect{AG}=\vect{AB}+\vect{BC}+\vect{CG}=\vi +\vj +\vec k$. De même, $\vect{BH}=\vect{BA}+\vect{AD}+\vect{DH}=-\vi +\vj +\vec k$.

\item \textit{[1 pt]} Par Chasles, $\vect{IJ}=\vect{AJ}-\vect{AI}=\frac{\vj -\vi }{5}=\frac{\vect{BD}}{5}$. Ces vecteurs sont non nuls et colinéaires : $(IJ)$ et $(BD)$ sont parallèles.

\item \textit{[1 pt]} Directeurs $\vi $ et $\vec k$, non colinéaires. Une intersection imposerait $t\vi =\vi +\vj +s\vec k$, impossible par le coefficient de $\vj $. Elles sont non coplanaires.

\item \textit{[1 pt]} $\vect{AO}=\frac{\vect{AG}}{2}$, donc $O\in(AG)\subset(ABG)$.

\item \textit{[1 pt]} Les points vérifient $\vect{IJ}=\frac{-\vi +\vj }{5}$ et $\vect{IK}=\frac{-\vi +\vec k}{5}$. Si $\vect{IK}=t\vect{IJ}$, le coefficient de $\vec k$ donnerait $\frac{1}{5}=0$, impossible. Les vecteurs ne sont pas colinéaires, donc $I,J,K$ ne sont pas alignés.

\end{enumerate}

\fi

\fi

